\section{问题重述}

数字娱乐行业中，策略类（SLG）游戏``代号K''即将开启新服运营。运营团队需要在保证用户留存（玩家持续活跃）的前提下，通过合理的资源投放与付费设计，最大化开服首月总收入。核心挑战在于留存率与付费率之间的最优平衡。

运营方提供了来自该游戏过往服务器、为期30天的真实玩家行为后台数据，包含2000名训练用户与1000名测试用户。实际数据由两类文件组成：资源变更记录（pickdata1, 33字段, $\sim$193万行）与用户行为事件记录（pickdata2, 52字段, $\sim$416万行），通过account\_id关联。

基于上述数据，需要逐步解决以下四个问题：

\textbf{问题1：玩家留存规律与流失预测模型。}计算并绘制次日、7日、14日、30日留存率曲线；识别玩家流失的关键时间节点；基于玩家前3天的行为数据构建数学模型，预测30天内的预期留存天数。

\textbf{问题2：资源消耗、成长速度与付费行为的关联分析。}量化资源消耗量与等级提升速度的关系，识别免费资源产出的``卡点''；分析钻石存量与流失风险的关系；分析首付特征及其对留存的影响；建立回归模型分析影响总付费金额的关键因子。

\textbf{问题3：新服首月付费策略设计与收益最大化模型。}将玩家群体划分为不同聚类并描述行为特征；假设新服导入10,000名行为同分布的新玩家，设计付费礼包投放与定价策略；在确保30日留存率$\geq 10\%$的约束下，最大化30天总营收（目标70,000元），并给出目标达成的概率分布证明。

\textbf{问题4：数据采集建议与策略闭环。}建议优先采集两类新埋点数据，详述理由，并阐述如何修正问题3中的定价与投放策略。
